% ===== Acknowledgements =====
\newpage
\goldrule
\begin{center}
{\large\bfseries\color{rosedeep} Acknowledgements}
\end{center}
\goldrule

\vspace{0.8em}

\noindent
{\cjk 感谢那一位纯粹、自然、坚韧、智慧的"森林与山之女孩"。}

\vspace{0.4em}

\noindent
\textit{My deepest gratitude to the one I call the girl of the forest and the mountains---pure, natural, steadfast, and wise. The concept of the happiness jar began, as all the best concepts do, not in a library but in a life: in the quality of her presence in the shared field, in the patient example of her attentiveness, in the way she showed me, without saying so, that the small contingent things are the most lasting. May we continue to fill the jar together.}

\vspace{1.2em}

\noindent
{\cjk 也愿世界上所有有趣的人，在反复的日常生活中，受于偶然，得之幸福。}

\vspace{0.4em}

\noindent
\textit{And may all the interesting people in this world---caught in the repetitions of ordinary life, in the daily texture of shared existence---find, through contingency, their wealth. Not the gold that can be counted, but the jar that grows fuller with every taking.}

\vspace{1.8em}

\goldrule

\vspace{1.2em}

\noindent{\color{rosedeep}\textbf{Use of Generative AI}}

\vspace{0.4em}

\noindent
The philosophical framework, theoretical claims, and original arguments of this paper are the author's own. The conceptual architecture---including the concept of Generative Relational Wealth, the structural remainder $\mathcal{R}(O(x), x) \neq 0$, the Relational Isomorphism Principle, the von Neumann problem for relational systems, the dialectical analysis of wealth in intimate relations, the epistemic justice framework applied to GRW, and the happiness jar as both theoretical paradigm and practical discipline---was developed through extended dialogue with Claude (Anthropic), a large language model, which served as an interactive writing and development partner throughout the drafting process.

Specifically, Claude assisted with: the drafting and \LaTeX\ typesetting of all sections; the identification and organisation of relevant literature across political economy, psychoanalysis, feminist theory, semiotics, and formal systems; the formal articulation of arguments whose substance had been established in prior discussion; and the production of the final compiled document.

The author reviewed, directed, and took intellectual responsibility for all content. The core philosophical positions, the selection of frameworks, all original claims, and all open questions were determined by the author. Claude did not independently generate theoretical positions; it developed and articulated positions established in collaborative dialogue.

\vspace{1.0em}

\goldrule
